\documentclass{article}


\usepackage{arxiv}

\usepackage[utf8]{inputenc} % allow utf-8 input
\usepackage[spanish]{babel}
\usepackage[T1]{fontenc}    % use 8-bit T1 fonts
\usepackage{hyperref}       % hyperlinks
\usepackage{url}            % simple URL typesetting
\usepackage{booktabs}       % professional-quality tables
\usepackage{amsfonts}       % blackboard math symbols
\usepackage{nicefrac}       % compact symbols for 1/2, etc.
\usepackage{microtype}      % microtypography
\usepackage{lipsum}
\usepackage{graphicx}
\usepackage{fancyhdr}
\usepackage{float}


\title{TP 1.1 Simulación de una ruleta}


\author{
 Tomás Lardizábal \\
  Legajo 47433\\
  Universidad Tecnológica Nacional - FRRO\\
  Zeballos 1341, S2000, Argentina \\
  \texttt{tlardizabal1@gmail.com} \\
  %% examples of more authors
   \And
 Iñaki Díaz\\
  Legajo 48944\\
  Universidad Tecnológica Nacional - FRRO\\
  Zeballos 1341, S2000, Argentina \\
  \texttt{inakidiaz2808@gmail.com} \\
    \And
 Tomás Splivalo\\
  Legajo 51665 \\
  Universidad Tecnológica Nacional - FRRO\\
  Zeballos 1341, S2000, Argentina \\
  \texttt{tomassplivalo1@gmail.com} \\
    \And
 Luciano Armas\\
  Legajo 47181\\
  Universidad Tecnológica Nacional - FRRO\\
  Zeballos 1341, S2000, Argentina \\
  \texttt{lucianoarmas11@gmail.com} \\
}

\date{\today}
\setlength{\itemsep}{0pt}
\begin{document}
\maketitle

\pagestyle{fancy}
\fancyhead{} % Elimina encabezado predeterminado
\rhead{Simulación de una ruleta, \today}

\begin{abstract}
Mediante el siguiente estudio se pretende realizar un análisis y evaluación de las distintas estrategias de juego en la ruleta a través de simulaciones en lenguaje Python, contemplando dos escenarios distintos: 
\begin{itemize}
  \item Uno en el cual el jugador posee un capital infinito
  \item Otro escenario en el que el jugador posee un capital acotado.
\end{itemize}
 El objetivo es conocer si hay posibilidades de lograr resultados positivos que hagan crecer el capital apostado.
Se evalúan las estrategias Martingala, Martingala Invertida, Fibonacci y D'Alembert.
\end{abstract}


\section{Introducción}
El juego de la ruleta ha sido analizado desde múltiples enfoques, tanto lúdicos como científicos, para responder a la pregunta de si existe algún algoritmo que pueda predecir la aleatoriedad de este juego, permitiéndonos obtener beneficios.
En este informe pondremos a prueba algunas de las famosas estrategias de apuestas para analizar su comportamiento a lo largo de múltiples tiradas y estudiar las probabilidades de obtener ganancia con cada estrategia, tomando una ruleta europea ideal.



\end{document}